\documentclass[a4paper,12pt]{article}
\usepackage[T2A]{fontenc}
\usepackage[utf8]{inputenc}
\usepackage[russian]{babel}
\usepackage{amsmath,amssymb,mathtools}
\usepackage{geometry}
\geometry{left=20mm,right=20mm,top=20mm,bottom=23mm}
\renewcommand{\familydefault}{\sfdefault}
\usepackage{microtype}
\usepackage{array,booktabs,tabularx}
\usepackage{tikz}
\usepackage{xcolor}
\usepackage{fancyhdr}
\usepackage[unicode=true,colorlinks=true,linkcolor=darkaccent]{hyperref}
\definecolor{accent}{HTML}{087C83}
\definecolor{darkaccent}{HTML}{07595D}
\definecolor{ink}{HTML}{27333B}
\definecolor{wrong}{HTML}{A53639}
\color{ink}
\setlength{\parindent}{0pt}
\setlength{\parskip}{5pt}
\linespread{1.13}
\pagestyle{fancy}
\fancyhf{}
\renewcommand{\headrulewidth}{0.25pt}
\fancyhead[L]{\scriptsize \textcolor{darkaccent}{МАТЕМАТИКА / РЕВЬЮ ДОМАШНЕЙ РАБОТЫ}}
\fancyhead[R]{\scriptsize 10.11.2026}
\fancyfoot[C]{\scriptsize Лёвин Артём Александрович эксклюзивно для Николь Саркисьянц}\fancyfoot[R]{\scriptsize\thepage}
\newcommand{\errorlabel}{\textcolor{wrong}{\textbf{Ошибка.}}}
\newcommand{\keyidea}{\textcolor{darkaccent}{\textbf{Ключевая идея.}}}
\newcommand{\result}{\textcolor{darkaccent}{\textbf{Итог.}}}
\newcommand{\subline}[1]{\par\vspace{3pt}{\color{darkaccent}\bfseries #1}\par}
\begin{document}
\thispagestyle{empty}
\begin{center}
\vspace*{28mm}
{\color{darkaccent}\large МАТЕМАТИКА}\par\vspace{15mm}
{\Huge Ревью домашней работы}\par\vspace{11mm}
{\fontsize{17}{23}\selectfont\color{accent}Тригонометрическое уравнение}\par\vspace{15mm}
{\small Подготовка к ЕГЭ по профильной математике}\par
\vfill
{\normalsize Николь Саркисьянц}\par\vspace{8mm}
{\normalsize 10 ноября 2026 г.}\par\vspace{12mm}
{\small Лёвин Артём Александрович}\par
{\small эксклюзивно для Николь Саркисьянц}\par\vspace{12mm}
\begin{tikzpicture}\draw[accent,line width=.9pt] (0,0) .. controls (1.1,.35) and (2.1,-.35) .. (3.2,0) .. controls (4.4,.35) and (5.7,-.35) .. (7.0,0);\end{tikzpicture}
\end{center}
\clearpage
\section*{Сводная таблица}
\textcolor{darkaccent}{Источник:} «Ревью ревью 09.11.26.pdf», упражнение к заданию 14.

В таблицу включены только задания, в которых по повторной проверке обнаружены ошибки. Номер задания ведёт к подробному разбору.
\par\par\vspace{4mm}
\small\renewcommand{\arraystretch}{1.35}
\begin{tabularx}{\linewidth}{@{}p{12mm}p{27mm}X p{32mm}p{35mm}@{}}
\toprule
\textbf{№} & \textbf{Тема} & \textbf{Суть ошибки} & \textbf{Ваш ответ} & \textbf{Верный ответ}\\\midrule
\hyperref[task:14]{14} & Тригонометрия & Пропущен корень и не проверен знак косинуса. & \shortstack[l]{$\frac\pi2,\ \frac{7\pi}6,$\\$\frac{11\pi}6$} & \shortstack[l]{$\frac\pi2,\ \frac{3\pi}2,$\\$\frac{11\pi}6$}\\
\bottomrule
\end{tabularx}
\normalsize
\clearpage
\thispagestyle{empty}
\vspace*{\fill}
\begin{center}{\fontsize{25}{30}\selectfont\color{darkaccent}Разбор ошибок}\par\vspace{5mm}
{\large Подробное объяснение и одна задача на закрепление}\end{center}
\vspace*{\fill}
\clearpage
\phantomsection\label{task:14}
\section*{Задание № 14}

\subline{Условие}
Решите уравнение и найдите все корни на отрезке $[0;2\pi]$:
\[
\sqrt{2\sin^3x-\cos^2x-2\sin x+\cos x}
=\sqrt{\sin\left(\frac\pi2-x\right)}.
\]
\subline{Ваше решение}
Правильно использовали $\sin(\frac\pi2-x)=\cos x$, приравняли подкоренные выражения и получили
\[
-\cos^2x(2\sin x+1)=0.
\]
После этого нашли семейства $\cos x=0$ и $\sin x=-\frac12$. Однако на заданном отрезке выписали только один корень первой ветви и оставили корень второй ветви, для которого правая часть исходного уравнения не определена.

\textbf{Ваш ответ:} $\frac\pi2$, $\frac{7\pi}6$, $\frac{11\pi}6$.

\subline{Где именно возникает ошибка}
\textbf{Последний правильный шаг:} $\cos x=0$ или $\sin x=-\frac12$.

\errorlabel\ При отборе на $[0;2\pi]$ пропущено $x=\frac{3\pi}2$ из первой ветви и ошибочно оставлено $x=\frac{7\pi}6$ из второй. В последней точке $\cos x<0$, поэтому $\sqrt{\cos x}$ в действительных числах не существует.

\subline{Почему этот шаг неверен}
Равенство двух арифметических квадратных корней требует, чтобы их подкоренные выражения совпадали и были неотрицательными. После перехода к $\sqrt{\cos x}$ обязательно сохраняем ограничение $\cos x\ge0$. При этом $\cos x=0$ допускается: корень из нуля определён. На отрезке $[0;2\pi]$ у уравнения $\cos x=0$ два решения: $\frac\pi2$ и $\frac{3\pi}2$.

\keyidea\ Сначала определите область допустимых значений корней, затем решите равенство подкоренных выражений и отдельно проверьте каждый найденный корень.

\subline{Исправление с последнего правильного шага}
\[
\begin{array}{ll}
\cos x=0&\Rightarrow\ x=\frac\pi2,\ \frac{3\pi}2;\\[3pt]
\sin x=-\frac12&\Rightarrow\ x=\frac{7\pi}6,\ \frac{11\pi}6.
\end{array}
\]
С учётом $\cos x\ge0$ исключаем $\frac{7\pi}6$. Остальные три корня сохраняются.

\subline{Полное правильное решение}
\textbf{1. Преобразуем правую часть.} По формуле приведения $\sin(\frac\pi2-x)=\cos x$. Поэтому исходное уравнение имеет вид
\[
\sqrt{A(x)}=\sqrt{\cos x},\quad A(x)=2\sin^3x-\cos^2x-2\sin x+\cos x.
\]
\textbf{2. Переходим к системе.} Из равенства корней следует $A(x)=\cos x$ и $\cos x\ge0$. Обратно, эти два условия обеспечивают равенство обоих определённых корней. Получаем
\[
\begin{cases}
2\sin^3x-\cos^2x-2\sin x=0,\\
\cos x\ge0.
\end{cases}
\]
\textbf{3. Упростим первое уравнение.} Используем $\sin^2x-1=-\cos^2x$:
\begin{align*}
2\sin x(\sin^2x-1)-\cos^2x&=0,\\
-2\sin x\cos^2x-\cos^2x&=0,\\
-\cos^2x(2\sin x+1)&=0.
\end{align*}
\textbf{4. Решим уравнение произведения.}
\[
\cos x=0\quad\text{или}\quad\sin x=-\frac12.
\]
На $[0;2\pi]$ первая ветвь даёт $\frac\pi2$ и $\frac{3\pi}2$, вторая --- $\frac{7\pi}6$ и $\frac{11\pi}6$.

\textbf{5. Применим ограничение.} При $x=\frac{7\pi}6$ косинус отрицателен. Для остальных точек косинус равен нулю или положителен.

\subline{Проверка}
При $x=\frac\pi2$ и $x=\frac{3\pi}2$ косинус равен нулю, а из найденного тождества следует $A(x)=0$. Обе части исходного уравнения равны нулю. При $x=\frac{11\pi}6$ имеем $\sin x=-\frac12$, $\cos x=\frac{\sqrt3}2>0$; после сокращения равенства подкоренных выражений получаем верное тождество, поэтому корни равны.

\result\ На указанном отрезке $\displaystyle x\in\left\{\frac\pi2,\frac{3\pi}2,\frac{11\pi}6\right\}$.

\subline{Задача на закрепление}
Решите уравнение и найдите все корни на отрезке $[0;2\pi]$:
\[
\sqrt{4\cos^3x+2\sin^2x-4\cos x+\sin x}
=\sqrt{\cos\left(\frac\pi2-x\right)}.
\]
\end{document}